\documentclass[11pt,letterpaper]{article}
\usepackage[top=0.6in, bottom=0.45in, left=0.7in, right=0.7in]{geometry}
\usepackage{lmodern}
\usepackage{hyperref}
\hypersetup{colorlinks=false, urlcolor=black, pdfborder=0 0 0}
\setlength{\parindent}{0pt}
\setlength{\parskip}{0.75em}
\raggedright
\pagestyle{empty}

\begin{document}

\noindent
{\bfseries Ke-Cheng Hsu (Clio)}\\
Ph.D.\ Candidate, Department of Economics, HKUST\\
(852) 6840\,9753 \quad $|$ \quad \href{mailto:khsu@connect.ust.hk}{khsu@connect.ust.hk}\\
\href{https://cliohsu1997.github.io/}{cliohsu1997.github.io} \quad $|$ \quad \href{https://github.com/cliohsu1997}{github.com/cliohsu1997}

\vspace{0.5em}
\noindent
16 November 2026

\vspace{0.35em}
\noindent
Hiring Team\\
UBS --- Research, Economics \& Strategy, China Equity Strategy, Shanghai\\
Re: Research -- Economics \& Strategy -- China Equity Strategy Associate (345218BR)

\vspace{0.2em}
Dear Hiring Team,

I am writing to apply for the Research -- Economics \& Strategy -- China Equity Strategy Associate position in Shanghai (345218BR). I am a Ph.D.\ candidate in Economics at HKUST (expected June 2027; I can finish as early as February 2027), a Hong Kong permanent resident, a native Mandarin speaker, and professionally fluent in English. I am drawn to the role because it combines careful analysis of China A-share markets---market developments, earnings, economic data, and investor flows---with clear writing for clients and colleagues. My training is in building frameworks that connect evidence to theory, estimating them on large data, and presenting the results so non-specialists can follow them.

In my job market paper on Taobao apparel return policy, I clean and link a large Chinese e-commerce transaction panel (2017--2023; \(\sim\)62M RMB revenue; six categories), estimate a structural model of purchase, return-shipping insurance, and returns, and compare policies by seller profit, consumer surplus, and social welfare. A counterfactual menu that pairs a discounted price with a moderate restocking fee raises social welfare by about 2 million RMB (roughly a 5\% gain) and seller profit by about 1.2 million RMB relative to the status-quo insurance policies. The work is close to what the posting describes day to day: maintain models and databases, run quantitative analysis, and turn findings into written notes under a clear research question.

With Professors Kohei Kawaguchi and Kei Kawai, I build title-by-market-by-month panels from bookstore point-of-sale records and estimate pricing and inventory under demand uncertainty. Relative to a classical static model, allowing for demand uncertainty and inventory lowers estimated marginal cost by about 20\%. We summarize results in figures, tables, and a weekly report on a fixed deadline. I have also presented related work at the Asia-Pacific Industrial Organization Conference, the Econometric Society Australasia Meeting, and the China Economics Annual Conference in 2026.

I run analysis pipelines in Python, R, and C++ on AWS and HPC/SLURM, and I use Excel, Word, and PowerPoint for models, figures, and presentations. I use AI tools regularly for research and coding, and I check outputs carefully before relying on them. Earlier work on an interbank lending network in MATLAB gave me some exposure to financial-market dynamics. As a lecturer and teaching assistant at HKUST (ratings 4.53/5 and 4.72/5), and as Teaching Assistant Coordinator in the Business School, I have practice explaining technical material clearly in Mandarin and English and keeping written work consistent before it goes out. I would welcome the chance to contribute to China equity strategy research in Shanghai.

I have attached my resume and transcript, and I can send my job market paper on request. Thank you for your consideration.

\vspace{0.5em}
Sincerely,\\[0.4em]
Ke-Cheng Hsu (Clio)

\end{document}
